\begin{textsample}{POS Dim 3 – persona\_gemini – Score 18.00 – t400\_gemini.txt}  \label{ex:f3_pos_045}
The \textbf{transport} \textbf{sector} is responsible for over a quarter of the EU 's \textbf{greenhouse} \textbf{gas} \textbf{emissions}.
In a stark failure of climate policy, these \textbf{emissions} have risen by 28 % since 1990.
Instead of reversing this \textbf{trend}, governments have \textbf{used} post-COVID recovery funds for \textbf{investments} in polluting carmakers and airlines, a move that prioritizes individualism over the societal \textbf{good}.
This path is unsustainable and requires immediate, radical change.
Governments must stop all \textbf{investments} in new airports and highways and \textbf{end} \textbf{tax} \textbf{breaks} that benefit \textbf{cars}.
Short-haul flights must be banned wherever a viable \textbf{alternative} exists, and \textbf{sales} of new combustion or hybrid \textbf{cars} must be halted by 2028.
The funds currently propping up polluting \textbf{industries} should be redirected to achieve decarbonization by 2040.
This \textbf{means} investing in affordable trains, \textbf{cycling} \textbf{infrastructure}, and \textbf{public} \textbf{transport}.
As we reshape our \textbf{transport} \textbf{systems}, it is essential to support worker transitions into green jobs. \textbf{Public} action is critical.
Individuals can contribute by contacting their representatives to demand change and by tracking their own personal \textbf{emissions} to understand their impact.

% matched lemmas: alternative, break, car, cycle, emission, end, gas, good, greenhouse, industry, infrastructure, investment, mean, public, sale, sector, system, tax, transport, trend, use
\end{textsample}
